%% notes-tex/publishing-system/sections/1-two-families.tex

\section{Two document families, one machinery\secstatus{todo}}
\label{sec:families}

The \term{manuscript} is the Nature Biotechnology paper in
\texttt{paper/nature-biotech/}. It uses the Springer Nature \texttt{sn-jnl} class
with the Nature Portfolio style, keeps its body in one file that three thin
wrappers include, and is the only document whose bibliography must come from the
group library alone. A \term{typeset note} is any document under
\texttt{notes-tex/<slug>/}: a design memo, a method review, an analysis report such
as the additive-baselines documents. It uses the \texttt{article} class with the
shared \texttt{tcdoc.sty}, which reproduces the manuscript's drafting look, status
chips and figure macros on top of a plain class that tolerates numbered sections
and a contents page.

\notesec{What they share}
Both are compiled by Tectonic, a single binary with no TeX Live install, which
fetches packages deterministically. Both export \texttt{SOURCE\_DATE\_EPOCH} as the
newest modification time among the build inputs, so hyperref writes no wall-clock
creation date and two builds of identical sources are byte-identical; that is what
lets a review copy be versioned by content hash. Both take a committed
\texttt{references.bib} that \texttt{make bib-pull} fetches from the served
bibliography store, verified against a manifest. Both publish through
\file{notes-tex/common/zotero_publish.py}, which derives the Zotero collection from
the directory path. And both obey the same rule about content: no figure, table or
number is authored inside a document; each is produced by a committed script under
\texttt{experiments/<id>/} and only converted, placed or checked by \texttt{make}.

\notesec{What differs}

%% SOURCE: hand-authored from paper/nature-biotech/Makefile and
%% notes-tex/common/Makefile.common at the commit this document was built from.
\begin{table}[htbp]
\centering
\footnotesize
\caption{The two families side by side.}
\label{tab:families}
\begin{tabular}{lp{0.36\textwidth}p{0.36\textwidth}}
\toprule
 & Manuscript & Typeset note \\
\midrule
Directory & \texttt{paper/nature-biotech/} & \texttt{notes-tex/<slug>/} \\
Makefile & its own, 250 lines & four lines, includes \texttt{common/Makefile.common} \\
Class & \texttt{sn-jnl}, Nature Portfolio style & \texttt{article} with \texttt{tcdoc.sty} \\
Body & \texttt{content.tex}, included by three wrappers & \texttt{<slug>.tex} root, \texttt{sections/*.tex} \\
Default build & \texttt{make paper}, three PDFs & \texttt{make}, one PDF named for the directory \\
Figure sources & draw.io, exported headless & script-written SVG panels, converted; or draw.io \\
Figure gate & hard: any view refuses to build while a figure is out of spec & \texttt{make check}, after the build \\
Bibliography & group \texttt{paper} collection only & union of group and personal collections of the same name \\
Publish & \texttt{make publish}, \texttt{editing.pdf} & \texttt{zotero\_publish.py <slug>}, the default PDF \\
Collaborators & \texttt{sync-overleaf.sh} to the shared Overleaf clone & the Zotero review copy \\
\bottomrule
\end{tabular}
\end{table}

\notesec{Vocabulary}
A \term{view} is one PDF rendering of the manuscript's shared body. The
\term{gate} is a check that reads a build's log and sources and exits non-zero on
a violation; the manuscript's runs before Tectonic and blocks the build, the note's
runs after it and blocks nothing but the landing. A \term{status chip} is the
colored mark after a heading, \texttt{todo}, \texttt{tent} or \texttt{final}, that
appears in the contents so the contents page is the status board. A
\term{provenance flag} is the visible marker \texttt{\textbackslash external} or
\texttt{\textbackslash secondhand} on a fact not backed by the library mirror; it
prints in every view. The \term{mirror} is the on-disk copy of the group's papers
under \texttt{\$DATA\_ROOT/torchcell-library/}, and \term{tc-lit} is the read-only
service that streams it and the bibliographies to any machine with a key.
